% !TeX root = ../main.tex

\chapter{Requirements en iteratief proces}
\label{chap:requirements}

\pending{Dit hoofdstuk inhoudelijk nakijken en afwerken. Controleer vooral de koppeling met het oorspronkelijke stageplan, beperk herhaling met de technische hoofdstukken en verplaats persoonlijke reflectie naar het reflectiehoofdstuk.}

\section{Iteratief overleg met de klant}
De requirements lagen bij de start niet volledig vast. Jan wist wel wat de belangrijkste informatie was, maar de precieze voorstelling was nog niet duidelijk. Daarom werd gedurende deze zes weken meerdere keren samen gezeten en overlegd over layout en nodige informatie zoals gemiddelde rondetijden, pitstopprojecties en sectorvergelijkingen. Voor mij was dit iteratief proces zeer leerrijk, doordat hoe ik sommige informatie voorstelde in het dashboard niet altijd duidelijk waren voor Jan. Tijdens dit proces heb ik dus geleerd om beter te luisteren naar wat anderen verwachten en nodig hebben. Achter de schermen was het ook voor mij een iteratief proces. De manieren van informatie verzamelen, verwerken en weergeven werden meerdere keren volledig aangepast. 

Na wat overleg en een eerste schets van het dashboard\ref{fig:dashboard-schets1}, werd er een eerste werkende versie gemaakt. Deze versie konden we samen testen tijdens een race in Spa-Francochamps. Tijdens deze race zagen we dat alle informatie rond rondetijden en sectoren correct werd weergegeven. Toch waren er een aantal dingen die beter moesten worden weergegeven of verzamelt en berekent. In een later hoofdstuk zal ik dieper ingaan op de problemen die we tegenkwamen tijdens deze race.~\ref{chap:spa}

\section{Evolutie van de requirements}
De eerste versie van het dashboard werd voornamelijk ontwikkeld als conceptuele demonstratie voor de eerste meeting. Deze versie gaf een eerste beeld van hoe de applicatie eruit zou kunnen zien en welke functionaliteiten technisch mogelijk waren. Ze bevatte onder meer een algemene sessiekaart, een klassentabel en enkele grafieken. 
Feedback leidde tot een volledig hertekend scherm volgens een schets~\ref{fig:dashboard-schets1} van de opdrachtgever. De informatie werd opgesplitst in grote tijdvakken, compacte sectorvakken, vergelijkingskolommen en een pitstopbalk. Zo kwamen we tot een eerste werkende versie van het dashboard~\ref{versie1}. Hier had ik ook nog een grafiekenvenster~\ref{fig:grafiekvenster1} toegevoegd om tijdens de race visueel de evolutie van de rondetijden te kunnen volgen.

Ook de achterliggende regels evolueerden. De oorspronkelijke referentietijdlogica werd eerst
verwijderd omdat de betekenis van de regel nog niet voldoende duidelijk was. Nadat de gewenste
werking was verduidelijkt, werd ze opnieuw toegevoegd als configureerbare lap- en sectorreferentie
met kleurstatussen. De pitstopprojectie begon als een eenvoudige aftrek van de pitstoptijd. Tests
toonden vervolgens dat rekening moest worden gehouden met alle tussenliggende wagens, ook uit andere
klassen, met rondeachterstanden, providerafhankelijke intervallen en FCY-omstandigheden.

De opslagstrategie veranderde ook. Eerst werd automatisch een nieuwe map per start gemaakt. Dat
was handig, maar onveilig wanneer de applicatie per ongeluk werd gesloten: een herstart zou een
nieuwe sessie beginnen. De uiteindelijke keuze is dat de gebruiker zelf een map per sessie
selecteert. Bij een herstart wordt bestaande JSONL-historiek geladen en worden alleen nieuwe ronden
toegevoegd waardoor de applicatie na een crash of sluiting gewoon verder kan gaan.

Dit waren slechts een aantal kleine voorbeelden van hoe verschillende manieren van opslag of berekeningen moesten veranderen om de juiste informatie te kunnen weergeven.

\section{Functionele vereisten}
De gerealiseerde functionele vereisten kunnen als volgt worden samengevat:

\begin{enumerate}
  \item De gebruiker configureert timing-URL, sessiemodus, opslagmap en aantal auto's die gevolgd moeten worden.
  \item De applicatie vraagt informatie van de timingpagina om de vijf seconden op en verdeelt dezelfde data over alle dashboards.
  \item De parser ondersteunt GetRaceResults en RIS Timing, deze twee timingpagina's zijn de enigste die gebruikt worden, en bewaart de bronprovider.
  \item Voltooide ronden worden persistent opgeslagen.
  \item Race-, practice- en qualifyingmodus tonen aangepaste analyses.
  \item Rondes onder FCY, Safety Car of andere neutralisatie blijven opgeslagen maar worden uit
        representatieve tempoanalyses gehaald.
  \item Inlaps en outlaps blijven opgeslagen maar tellen niet mee voor tempovergelijkingen.
  \item Groene sectoren uit een gedeeltelijk geneutraliseerde ronde kunnen afzonderlijk geldig
        blijven, terwijl de volledige ronde ongeldig is voor het rondetijdgemiddelde.
  \item De voorspelde rondetijd gebruikt actuele sectoren en recente sectordata van de huidige
        piloot.
  \item De gebruiker kan lap- en sectorreferenties aanpassen en ziet de veiligheidsmarge via kleuren
        en delta's.
  \item De pitstopplanner houdt rekening met openingsvensters, cooldown, verplichte stops en een
        instelbare pitstoptijd.
  \item Grafieken tonen piloot-, sector- en klassevergelijkingen en ondersteunen zoomen bij lange
        reeksen.
  \item De app kan na een crash of sluiting verdergaan vanuit dezelfde sessiemap.
\end{enumerate}

\section{Niet-functionele vereisten}
Tijdens een race wegen betrouwbaarheid en begrijpelijkheid zwaarder dan theoretische complexiteit.
De berekeningen moesten daarom deterministisch, uitlegbaar en onafhankelijk van de renderer zijn. De
renderer krijgt kant-en-klare viewmodellen en formatteert of berekent geen racebeslissingen zelf.
Elke belangrijke domeinfunctie moest met synthetische en historische data testbaar zijn.

Verder moest de applicatie langdurig kunnen draaien zonder onbegrensde geheugengroei of overmatige
opslag. De pollingfrequentie van vijf seconden levert tijdens 24 uur 17.280 meetmomenten op, maar
alleen gewijzigde snapshots en voltooide rondes worden als relevante historiek bewaard. Het gebruik
van één collector voor meerdere dashboards voorkomt onnodige netwerkbelasting.


\iffalse
\section{Traceerbaarheid van de requirements}

Tabel~\ref{tab:requirements-traceability} verbindt de belangrijkste requirements met hun implementatie, validatie en de correcties die tijdens de ontwikkeling werden uitgevoerd. Hierdoor wordt zichtbaar hoe feedback uit tests en praktijksituaties naar concrete verbeteringen werd vertaald.

\small
\begin{longtable}{@{}
  >{\raggedright\arraybackslash}p{2.8cm}
  >{\raggedright\arraybackslash}p{3.5cm}
  >{\raggedright\arraybackslash}p{3.5cm}
  >{\raggedright\arraybackslash}p{3.2cm}
@{}}
\caption{Traceerbaarheid van requirements, implementatie en evaluatie.}
\label{tab:requirements-traceability} \\

\toprule
\textbf{Requirement} &
\textbf{Implementatie} &
\textbf{Test en evaluatie} &
\textbf{Belangrijkste correcties} \\
\midrule
\endfirsthead

\toprule
\textbf{Requirement} &
\textbf{Implementatie} &
\textbf{Test en evaluatie} &
\textbf{Belangrijkste correcties} \\
\midrule
\endhead

\bottomrule
\endfoot

Live timingdata ophalen &
Een verborgen Electron-venster leest periodiek de timingpagina. Providerspecifieke parsers ondersteunen RIS Timing en GetRaceResults. &
Parsertests met opgeslagen HTML- en timingsnapshots, aangevuld met live praktijktests tijdens raceweekends. &
Ondersteuning toegevoegd voor afwijkende kolomnamen, ontbrekende kolommen en wijzigingen in de structuur van timingpagina's. \\

\addlinespace
Uniform gegevensmodel &
Providerspecifieke waarden worden omgezet naar een gemeenschappelijk model voor wagens, piloten, rondes, sectoren, gaps en sessiestatussen. &
Dezelfde berekeningsfuncties werden getest met gegevens van beide timingproviders. &
Interpretatie van \texttt{GAP}, \texttt{DIFF} en \texttt{INT} werd aangepast voor tijdsverschillen en rondeachterstanden. \\

\addlinespace
Lokale opslag en herstel &
Snapshots, rondehistoriek, sessiemetadata en berekende statistieken worden in leesbare JSON-, JSONL- en CSV-bestanden opgeslagen. &
Herstarttests met bestaande sessiemappen en controle op dubbele of ontbrekende ronden. &
Sessiemappen kunnen opnieuw worden geopend, zodat een race na een crash of onbedoelde afsluiting kan worden hervat. \\

\addlinespace
Alleen geldige tijden gebruiken &
Ronden en sectoren krijgen een status voor onder andere FCY, Safety Car, rode vlag, pit-in, outlap en track limits. &
Unit tests en controle met opgeslagen racedata waarbij gemiddelden handmatig werden nagerekend. &
Handmatige correctie van rondestatussen toegevoegd. De eerste raceronde, pitgerelateerde ronden en geneutraliseerde ronden worden uitgesloten. \\

\addlinespace
Ronde- en sectoranalyse &
De applicatie berekent beste tijden, gemiddelden, laatste tien geldige ronden, ideale rondetijd en voorspelde rondetijd. &
Tests met vaste datasets voor meerdere piloten, stints en weersomstandigheden. &
Berekeningen werden modulair gemaakt en aangepast zodat dry, wet, intermediate en gecombineerde data afzonderlijk kunnen worden bekeken. \\

\addlinespace
Vergelijking binnen de klasse &
Het dashboard vergelijkt de gevolgde wagen met teamgenoten, de beste wagen in de klasse en andere klassewagens. &
Testscenario's met meerdere wagens, piloten en wagens op verschillende ronden. &
Deltatekens en kleuren werden gecorrigeerd. Klassewagens worden volgens hun positie gerangschikt en rondeachterstanden worden afzonderlijk behandeld. \\

\addlinespace
Gap- en catchvoorspelling &
Actuele gaps worden gecombineerd met recente geldige rondetijden om te schatten of een wagen wordt ingehaald. &
Tests voor tijdsgaps, lap gaps, ontbrekende gaps en wagens uit andere klassen. &
Providerspecifieke interpretatie toegevoegd voor gaps naar de leider en intervallen tussen opeenvolgende wagens. \\

\addlinespace
Pitstopplanning &
De planner bewaakt verplichte stops, gesloten pitperiodes, minimale tijd tussen stops, verwachte pit loss en de laatst veilige pitstop. &
Unit tests voor grenswaarden en praktijkevaluatie met opgeslagen racedata. &
Pitstopprojecties houden rekening met rondeachterstanden, FCY en de werkelijk gemeten duur van eerdere pitstops. \\

\addlinespace
Meerdere wagens volgen &
Eén datacollector voedt maximaal drie afzonderlijke dashboardvensters met dezelfde timingsnapshot. &
Handmatige tests met meerdere gevolgde wagens en afzonderlijke instellingen per venster. &
Dataverzameling werd gedeeld om meerdere downloads van dezelfde timingpagina te vermijden. \\

\addlinespace
Automatische rapportering &
Na een stint of sessie worden PDF-verslagen gegenereerd met ronde-, sector-, piloot- en pitstopanalyse. &
Controle met de opgeslagen racedata uit Spa en vergelijking met handmatig berekende resultaten. &
Stintnummering, conditionele gemiddelden, pitstopoverzichten en rapportgeneratie bij handmatig of automatisch stoppen werden gecorrigeerd. \\

\addlinespace
Duidelijke en robuuste bediening &
Het dashboard toont de belangrijkste informatie op één scherm en ondersteunt race-, kwalificatie- en trainingsmodi. &
Iteratieve beoordeling door de opdrachtgever en praktijktests op macOS en Windows. &
De UI werd meerdere keren herschikt. Bevestiging bij stoppen, automatische sessiedetectie, dark mode en herstelbare instellingen werden toegevoegd. \\

\end{longtable}
\normalsize

\section{Planning over zes weken}
De uitvoering verliep niet als zes strikt gescheiden fasen. In de eerste periode lag de nadruk op
parser, live collector en opslag. Daarna volgden dashboardanalyse, tests en de nieuwe UI.
Pitstoplogica, voorspelling, sessiemodi en grafieken werden in opeenvolgende iteraties toegevoegd.
In de laatste fase kwamen provideredgecases, FCY-projecties, distributie, automatische updates en
herstel na herstart aan bod.

Deze iteratieve volgorde week af van een klassiek vooraf vastgelegd stageplan, maar paste beter bij
de beschikbaarheid van live demo's en klantfeedback. Bugs die alleen in een bijna volledige race
zichtbaar werden, hadden een grotere impact op de prioriteiten dan oorspronkelijk geplande
uitbreidingen. Daarom was de race in Spa-Francochamps een belangrijk moment om de voortgang te evalueren en de resterende tijd te plannen aangezien dit ook de enige race was voor de \emph{24h of Zolder}. Daarom heb ik tijdens een andere race op het circuit van Zolder een Duits team dat op dit evenement aanwezig was of ik samen met hun mijn dashboard extra zou kunnen testen. Ook hier zijn weer kleine aanpassingen uit voortgevloeid, vooral UI gerelateerde aanpassingen, maar alle vitale functionaliteiten werkten hier al perfect.
\fi